% ---------------------------------------------------------------------
%  IŠVADOS — dvi pastraipos.
%  Skaičiai čia NEPERRAŠOMI ranka iš atminties — visi jau yra
%  atitinkamų skyrių lentelėse.
%  2026-09-15: sutrauktos iš šešių punktų (po vieną kiekvienam
%  uždaviniui) į dvi pastraipas vadovo nurodymu. Ankstesnė redakcija —
%  git istorijoje.
% ---------------------------------------------------------------------

Darbo tikslas buvo sukurti ir eksperimentiškai įvertinti dirbtiniu
intelektu pagrįstą IoT tinklo kibernetinių atakų aptikimo sprendimą.
Tikslas pasiektas. Iš 18 apžvelgtų metodų dviem pakopomis parinkti
keturi, sukurta visa grandinė nuo žalio srauto iki pavojaus signalo ir
įvertinta nepriklausoma testavimo aibe. Geriausias variantas ties
reikalaujamu 1~\% klaidingų teigiamų biudžetu pasiekia macro-F1
\num{0.6626} prie 0,95~\% klaidingų teigiamų, aptinka 88,5~\% atakų
srauto, telpa į 45~MB ir vienam įrašui sunaudoja apie 30 mikrosekundžių,
t.~y. tris eiles mažiau, nei leidžia kraštinio šliuzo biudžetas.

Ribojantis veiksnys buvo ne metodai, o duomenų struktūra. Keturi iš
aštuonių metodų atmetimų kilo iš to, kad rinkinyje nėra nei laiko eilės,
nei srauto krypties, nei šaltinio identifikatorių, o bendras tikslumas
šiam uždaviniui netinka, nes 78,7~\% klaidų yra painiava tarp pačių
atakų. Du rezultatai galioja plačiau už šį rinkinį. Nugalėtojas
priklauso nuo sprendimo taško, nes ties didžiausios tikimybės tašku
pirmauja Random~Forest, o ties biudžetu, kuriuo sistema realiai dirbtų,
XGBoost; darbuose, skelbiančiuose tik pirmąjį skaičių, rikiuotė gali
būti kita. Paneigta ir darbo prielaida, kad nematytoms atakoms būtinas
neprižiūrimas metodas, nes apibendrinimo geba priklauso ne nuo
paradigmos, o nuo to, ar mokymo aibėje lieka gimininga klasė.
